% =============================================================
%  Symbols
%  Collects the symbols of Sections 4 to 9. Compound subscripts use
%  an underscore (U_{dpe\_rec}); complex quantities are underlined.
% =============================================================
\section{Symbols}
\label{sec:symbols}

AC voltages are line-to-line RMS values unless stated otherwise; DC voltages are pole-to-ground. Complex quantities
are underlined. Powers and currents carry the sign of the convention of Section~\ref{sec:conventions}.

\begin{longtable}{@{}l p{0.80\linewidth}@{}}
  \toprule
  \multicolumn{1}{c}{\textbf{Symbol}} & \textbf{Description} \\
  \midrule
  \endhead
  $a$        & Inner modulation index of the converter \\
  $B$        & Shunt susceptance of the AC line \\
  $C$        & Capacitance; $C'$ per unit length of the AC line, $C_{sm}$ of one submodule \\
  $c_p$      & Specific heat capacity of the cooling water \\
  $e$, $\underline{E}$ & Internal converter voltage (phase) behind the arm reactors; $e_3$ the third harmonic \\
  $E_{on}$, $E_{off}$, $E_{rec}$ & Turn-on, turn-off and reverse-recovery energy per pulse \\
  $f$        & Frequency; $f_r$ rated frequency, $f_{sw}$ average switching frequency of a submodule \\
  $I$        & Current; $I_d$ DC current of a pole, $I_{arm}$ arm current, $I_C$ rated collector current \\
  $i_0$      & No-load current of the transformer \\
  $k$        & DC voltage across one converter in multiples of $U_{dpe}$ \\
  $k_3$      & Peak factor of the arm voltage, 1 or $\sqrt{3}/2$ with third harmonic injection \\
  $L$        & Inductance; $L_{arm}$ of one arm reactor \\
  $l$        & Length of a line \\
  $m$        & Modulation index of the converter \\
  $N$        & Submodules per arm \\
  $n$        & Voltage ratio of the transformer; $n_{conv}$ converters per station, $n_{par}$ modules in parallel,
               $n_{ins}$ inserted submodules \\
  $P$, $Q$, $\underline{S}$ & Active, reactive and complex power \\
  $P_{cu}$, $P_{fe}$ & Load and no-load losses of the transformer \\
  $Q_w$      & Cooling water flow per cold plate \\
  $R$        & Resistance; $R'$ per unit length \\
  $R_{th}$   & Thermal resistance; $R_{th\_jc}$ junction to case, $R_{th\_cs}$ case to heat sink, $R_{th\_hw}$ heat
               sink to water \\
  $r$        & Slope resistance of the on-state model \\
  $T$        & Temperature; $T_j$ junction, $T_w$ water \\
  $u_k$      & Short-circuit voltage of the transformer \\
  $U$        & Voltage; $U_{dpe}$ DC voltage pole-to-ground at the station DC terminal, $U_{dc}$ across the
               converter, $U_{conv}$ internal converter voltage \\
  $u_{sm}$, $u_c$ & Average and instantaneous capacitor voltage of a submodule \\
  $v_0$      & Threshold voltage of the on-state model \\
  $W$        & Energy stored in the submodule capacitors of an arm; $W_0$ its average \\
  $X$        & Reactance \\
  $Z$        & Impedance; $Z_{th}$ transient thermal impedance \\
  $\alpha$   & Temperature coefficient of the conductor resistance \\
  $\rho$     & Density of the cooling water \\
  $\tau$     & Time constant of a Foster element \\
  $\omega$   & Angular frequency, $2 \pi f$ \\
  \bottomrule
\end{longtable}
